Merkur hat die grosse Halbachse $a = \aMeO$ und die siderische Umlaufzeit $T = \TMeO$ (FoTaBe S.~212). Auf einer Kreisbahn mit Radius $a$ legt er in der Zeit $T$ den Umfang $2\pi a$ zurück:
\begin{align}
 v &= \vMeF \\
   &= \frac{2\pi\cdot\aMe}{\TMe} \\
   &= \vMe \approx \result{\vMeP{3}{3}}
\end{align}
Die Abweichung vom tatsächlichen Mittelwert $v_\text{m} = \vReO$ ist
\begin{align}
 \frac{v-v_\text{m}}{v_\text{m}} &= \frac{\vMe-\vRe}{\vRe} = \dev \approx \result{\devP{0}{2}}
\end{align}
Der Kreis mit Radius $a$ ist etwas länger als die Ellipsenbahn: Bei der numerischen Exzentrizität $\varepsilon = \exMO$ ist der Umfang der Ellipse näherungsweise $2\pi a\left(1-\frac{\varepsilon^2}{4}\right)$. Damit erhält man $v = \vElF = \vElP{3}{4}$, in guter Übereinstimmung mit dem tatsächlichen Mittelwert.
